%% Section 10: Conclusion
\section{Conclusion}
\label{sec:conclusion}
The leakage current of H4 and HERIC transformerless inverters with an L filter was investigated as a function of switching frequency (1--40~kHz) by simulation. In all four cases examined the leakage current exhibits sharp peaks, up to 0.75~A (H4 unipolar), 0.5~A (H4 bipolar) and 4~A (HERIC unipolar), located at the CM resonance $\fres=1/(2\pi\sqrt{\Leq\Cpv})$ at its low-order sub-multiples or, in one case, at twice its value, as predicted from the simulation parameters to within about 5\% for the dominant peaks. Away from the resonance, the leakage falls to 15.4~mA (H4 unipolar, 8.5~kHz), 14~mA (H4 bipolar, 40~kHz), 40.3~mA (HERIC unipolar, 26~kHz) and 3.5~mA (HERIC bipolar, 13.5~kHz). For identical parameters, the HERIC bipolar configuration yields about eleven times lower leakage than HERIC unipolar, with a floor that is consistent with the line-frequency term $v_o/2$ of the effective CM source acting on $\Cpv$, but a low-order THD of about 9.1\% that is independent of $\fsw$ and attributable to zero-crossing distortion. The practical consequence is a design rule: the switching frequency, its sub-multiples and twice $\fres$ should be placed outside a $\pm20\%$ band, and $\fsw>2.4\fres$ is preferable.

Future work comprises: (i) re-simulation of all topologies (including H5 and H6) with identical $\Vdc$, $L$, $R_L$ and power; (ii) inclusion of switching and thermal loss models and of the ground-path resistance; (iii) a grid-connected, closed-loop model with grid impedance; (iv) LC and LCL filters; (v) a correction of the HERIC bipolar gating; and (vi) a hardware prototype to validate the resonance-based frequency-selection rule for several values of $\Cpv$.
